\chapter{Bone Marrow Transplant}

\section*{Overview}
A bone marrow transplant is a type of hematopoietic cell transplant (HCT), also called a blood-forming stem cell transplant. The name is commonly used even when the cells come from the bloodstream or donated umbilical cord blood rather than marrow. These cells can develop into red blood cells, white blood cells, and platelets.

The transplant may use the patient's own cells (autologous) or cells from a donor (allogeneic). A rare third type uses cells from an identical twin (syngeneic). The transplant team selects the cell source and approach for the person's disease and circumstances.

\section*{Common Instruments and Equipment}
The exact setup varies by technique and setting. Common items may include:
\begin{itemize}[nosep,leftmargin=*]
  \item central venous access catheter and infusion tubing
  \item blood-product and cell-infusion equipment
  \item cell-processing and cryopreservation equipment when applicable
\end{itemize}

\section*{Alternative Names}
Hematopoietic cell transplant (HCT); hematopoietic stem cell transplant (HSCT); blood or bone marrow transplant; BMT; marrow transplant.

\section*{Principle}
Before an allogeneic transplant, treatment called conditioning is given. It may include chemotherapy and sometimes radiation. Depending on the disease and treatment plan, conditioning may destroy diseased cells, suppress the immune system so donor cells can take hold, or do both. It is not always intended to completely destroy the marrow. Autologous transplants also use intensive treatment to treat the disease before the patient's stored cells are returned.

After infusion into a vein, the blood-forming cells travel to the marrow and begin making blood cells. This process is called engraftment. Engraftment takes time, and blood counts may remain low during the early recovery period.

\section*{History}
Bone marrow was the first widely used source of blood-forming stem cells for transplantation. Today, cells may also be collected from circulating blood or donated umbilical cord blood.

\section*{Why It Is Done}
HCT may be considered for selected blood cancers and serious diseases of the blood, marrow, or immune system. Examples include some leukemias, lymphomas, multiple myeloma, myelodysplastic syndromes, aplastic anemia, sickle cell disease, and certain inherited immune or metabolic disorders. The possible benefit depends on the specific disease and its status, other treatment options, transplant type, donor or cell source, and the person's overall health. A transplant is not appropriate for every person with these conditions and does not guarantee a cure.

\section*{Role in Care}
HCT is a treatment, not a diagnostic test. It may be a potentially curative treatment for selected conditions, or it may aim to control disease or improve long-term blood-cell production. The transplant specialist discusses the likely benefits, alternatives, and risks for the individual.

\section*{How It Is Performed}
The transplant team first identifies the cell source:
\begin{itemize}
  \item \textbf{Autologous:} The patient's own blood-forming cells are collected, usually from the bloodstream, and stored for later use.
  \item \textbf{Allogeneic:} Cells come from a related or unrelated donor, or from donated umbilical cord blood. Donor choice depends on tissue-type compatibility and other factors. A fully matched donor is not always required; a partially matched donor or donated cord blood may be an option.
  \item \textbf{Cell source:} Cells may be collected from the bloodstream, from marrow in the pelvic bones, or from cord blood. For a marrow collection, the donor is under anesthesia while marrow is withdrawn through needles placed in the back of the pelvic bones. To collect peripheral blood stem cells, the patient or donor usually receives growth-factor medicine for several days to move more stem cells into the blood. A machine then collects the cells from the blood and returns the remaining blood; this process is called apheresis and is not the same as surgical marrow collection.
\end{itemize}

The recipient receives a planned conditioning regimen. The cells are then infused through a central venous catheter, much like a blood transfusion; infusion usually does not require surgery or general anesthesia. The team monitors blood counts and clinical signs while the cells settle in and engraft.

\section*{Preparation}
Preparation is individualized and may include disease assessment, blood tests, heart, lung, kidney, and liver evaluation, infection screening, donor matching, and a review of medicines and other health conditions. The team explains the planned conditioning, cell source, expected hospital or outpatient care, and possible complications. A central venous catheter is commonly placed for treatment, transfusions, and blood sampling.

Because conditioning can affect fertility, patients should be offered a timely discussion of fertility-preservation options before treatment when appropriate. The care team also helps plan practical support during the period of close monitoring and recovery.

\section*{Contraindications}
There is no single list of conditions that automatically rules out every transplant. Do not stop prescribed medicines or delay reporting an infection while awaiting transplant without advice from the transplant team.

\section*{Precautions}
Significant heart, lung, liver, or kidney problems, an uncontrolled infection, the disease's status, and other health or support factors can increase risk or change the treatment plan. A transplant center assesses these factors and weighs them against the expected benefit.

\section*{Risks}
Risks vary with the disease, conditioning, cell source, and transplant type. They can include:
\begin{itemize}
  \item Low blood counts, which can lead to infection, anemia, and bleeding, and may require transfusions.
  \item Side effects of conditioning, such as mouth and throat sores, nausea, vomiting, diarrhea, fatigue, and injury to organs. A serious but uncommon liver complication is sinusoidal obstruction syndrome, formerly called veno-occlusive disease.
  \item Delayed engraftment or graft failure, relapse of the original disease, blood clots, and infertility or reduced fertility.
  \item \textbf{Graft-versus-host disease (GVHD):} This occurs mainly after an allogeneic transplant, when donor immune cells attack the recipient's tissues. It may affect the skin, gut, liver, eyes, mouth, lungs, or other organs. Rarely, GVHD has also been reported after autologous transplantation.
  \item Longer-term effects, including prolonged immune weakness, later infections, and a higher risk of some other cancers.
\end{itemize}

\section*{Aftercare and Recovery}
Care may take place in hospital, as an outpatient, or in both settings. The team checks blood counts and organ function regularly, watches for infection and other complications, and provides medicines, fluids, nutrition support, and red-cell or platelet transfusions when needed. Growth factors may be used in some treatment plans. People with an allogeneic transplant often take medicines to prevent or treat GVHD and need close monitoring for their effects.

The time to engraftment and discharge varies. Immune recovery can take much longer than the initial rise in blood counts, so infection precautions and frequent visits may continue for months or longer. Follow the transplant center's instructions for temperature checks, food and activity precautions, medicines, and when to call.

\section*{Outcome and Follow-up}
The team follows blood counts to assess blood-cell recovery. After an allogeneic transplant, testing for donor-cell chimerism (the proportion of blood-forming cells that come from the donor rather than the recipient) may help assess how the transplant is working; whether and how often it is checked depends on the clinical situation. Engraftment means blood-forming cells are working; by itself, it does not mean the original disease is in remission or cured.

Long-term follow-up checks for disease recurrence, GVHD after an allogeneic transplant, infections and immune recovery, organ or endocrine effects, fertility concerns, and late complications. People are usually revaccinated after HCT because protection from vaccines received before transplant may be lost. The transplant team sets the timing and schedule. Outcomes differ widely, and the team should explain the person's expected benefits and risks.

\section*{When to Contact the Care Team}
After transplant, contact the transplant team promptly for a fever or chills, new cough or sore throat, shortness of breath, bleeding, persistent vomiting or diarrhea, a new rash, yellowing of the skin or eyes, or another sudden change. Follow the center's specific fever threshold and emergency instructions. Seek emergency help for severe breathing difficulty, chest pain, fainting, or heavy bleeding.

\section*{References}
\begin{enumerate}
  \item National Health Service (NHS). What is a stem cell or bone marrow transplant? \url{https://www.nhs.uk/tests-and-treatments/stem-cell-or-bone-marrow-transplants/what-it-is/} (accessed November 23, 2026).
  \item National Health Service (NHS). How a stem cell or bone marrow transplant is done. \url{https://www.nhs.uk/tests-and-treatments/stem-cell-or-bone-marrow-transplants/how-its-done/} (accessed November 23, 2026).
  \item National Health Service (NHS). Complications of a stem cell or bone marrow transplant. \url{https://www.nhs.uk/tests-and-treatments/stem-cell-or-bone-marrow-transplants/complications/} (accessed November 23, 2026).
  \item NMDP. Patient eligibility for hematopoietic cell transplantation. \url{https://network.nmdp.org/services-support/hematology-oncology/consultation-guidelines-outcomes/patient-eligibility-for-hct} (accessed November 23, 2026).
  \item Centers for Disease Control and Prevention. Altered Immunocompetence: Recipients of Hematopoietic Cell Transplants. \url{https://www.cdc.gov/vaccines/hcp/imz-best-practices/altered-immunocompetence.html} (accessed November 23, 2026).
  \item NMDP. Understanding the Blood Stem Cell Transplant Process. \url{https://www.nmdp.org/patients/understanding-transplant/transplant-process} (accessed November 23, 2026).
  \item Kaya SY, et al. Myth or reality: Graft versus host disease after autologous hematopoietic cell transplantation, a multicentre experience. \textit{Transplant Immunology}. 2025;93:102286. \url{https://doi.org/10.1016/j.trim.2025.102286} (accessed November 23, 2026).
  \item European Society for Blood and Marrow Transplantation. Hepatic Complications. In: The EBMT Handbook: Hematopoietic Cell Transplantation and Cellular Therapies. 8th ed. 2024. \url{https://www.ncbi.nlm.nih.gov/books/NBK608328/} (accessed November 23, 2026).
  \item NMDP. How Bone Marrow and Blood Stem Cell Donors Are Matched. \url{https://www.nmdp.org/get-involved/join-the-registry/matching-with-a-patient} (accessed November 23, 2026).
\end{enumerate}

\clearpage
